% =====================================================================
%  Chapter 24 review -- the 14-station revision contract.
%  Revision-only: tables, one-liners, pointers back via \secref/\thmref/
%  \defref/\figref/\coderef/\probref.  No re-derivations.
%  Summit: CF 757G "Can Bash Save the Day?" (3400, L6) -- a weighted tree plus a
%  permutation a[1..n].  Query 1 l r v = sum over i in [l,r] of dist(a_i, v);
%  query 2 x = swap a_x and a_{x+1}; queries are XOR-encoded against the last
%  answer (online).  Centroid decomposition with a Fenwick / segment tree per
%  centroid over POSITIONS, storing dist(a_i, centroid), plus the child-side
%  "sub" subtraction so each pair is claimed at its top shared centroid.  Both
%  official samples hand-replayed (XOR-decoded): answers 23, 37, 28.
%  Exports prob:c18-summit, lst:c18-kit, lst:c18-summit.
%  Mixed set is EXACTLY 10, all live-verified, un-owned, sorted ASCENDING by
%  effective difficulty; each hint names its reduction and cites a c18 listing.
%  CF 342E (Ch.12), CF 161D (Ch.13), CF 914D/1601C/1743F and the whole
%  Ch.19--23 rosters are OWNED elsewhere -- never re-linked here.
% =====================================================================

\chaptersection{Chapter 24 Summary}

One sentence: \emph{a rooted DP can only see paths that pass through its own
root, so it is blind to the pairs that cross it; centroid decomposition fixes the
blindness by splitting every tree at its \emph{balance point}, claiming each path
at the single highest centroid it passes through, solving the pieces independently,
and merging a per-component distance table in \(\bigO{\log n}\) levels
(\(\secref{c18-tree}\), \(\secref{c18-count}\)).} Chapter~12 (\chref{c09}) proved a
centroid exists and finds one in linear time; this chapter starts at ``remove it
and recurse'' and ends with a live index and a product fold. Each row below is a
pointer, not a proof.

\begin{center}
\renewcommand{\arraystretch}{1.25}
\begin{tabular}{@{}p{0.24\linewidth} p{0.48\linewidth} p{0.20\linewidth}@{}}
\toprule
\textbf{Result} & \textbf{Statement} & \textbf{Where} \\
\midrule
Decomposition & Remove the centroid, mark it dead, recurse on each surviving piece; every
piece has \(\le n/2\) nodes so there are \(\bigO{\log n}\) levels and \(\bigO{n\log n}\) total
work. & \defref{c18-decomp}, \coderef{c18-build} \\
Centroid tree height & Parent \(\mathit{par\_c}\) links each centroid to the one that split
its piece; the tree has height \(\le \lceil\log_2 n\rceil\), so a node touches \(\bigO{\log n}\)
centroids. & \thmref{c18-height}, \figref{c18-centtree} \\
Top-centroid claim & Every simple path is claimed by exactly one centroid -- the highest one it
passes through -- so counting over all paths reduces to a per-centroid merge plus recursion. & \defref{c18-claim}, \secref{c18-count} \\
Distance histogram count & Collect each arm's depths, query-then-merge a histogram, and
\emph{subtract} the same-arm pairs; \(\bigO{n\log n}\) for all pairs at distance \(K\)
(\(\le K\)). & \coderef{c18-count}, \coderef{c18-715c} \\
Centroid distance index & One distance table per centroid; a node lives in its \(\bigO{\log n}\)
centroid-ancestor tables, so mark / nearest / within-budget queries are \(\bigO{\log n}\). & \defref{c18-index}, \thmref{c18-nn} \\
Beyond paths & Feature (additive along the arm), accept (a dominance / threshold test), and fold
(accumulator) are the only plugs that change from counting to products. & \defref{c18-beyond}, \coderef{c18-833d} \\
Engine map & Path \emph{ranges} are HLD, subtree \emph{sets} are sack, per-node \emph{states}
are DP; centroid owns \emph{all paths} and \emph{distance-to-set}. & \defref{c18-map}, \secref{c18-boundary} \\
\bottomrule
\end{tabular}
\end{center}

\medskip
{\small\textbf{Algorithms from this chapter.}}
\begin{center}
\renewcommand{\arraystretch}{1.2}
\begin{tabular}{@{}p{0.36\linewidth} l l >{\raggedright\arraybackslash}p{0.20\linewidth}@{}}
\toprule
\textbf{Algorithm} & \textbf{Prep} & \textbf{Per op} & \textbf{Code} \\
\midrule
Build centroid tree (find, mark dead, recurse) & \(\bigO{n\log n}\) & -- & \coderef{c18-build} \\
Count pairs at distance \(=K\) / \(\le K\) & \(\bigO{n\log n}\) & one pass & \coderef{c18-count} \\
CF 715C (Digit Tree) -- decimal-hash pair count & \(\bigO{n\log n}\) & one pass & \coderef{c18-715c} \\
Centroid distance index (mark / nearest) & \(\bigO{n\log n}\) & \(\bigO{\log n}\) & \coderef{c18-index} \\
CF 293E (Close Vertices) -- 2-D dominance count & \(\bigO{n\log^2 n}\) & one pass & \coderef{c18-293e} \\
CF 833D (Red-Black Cobweb) -- product fold & \(\bigO{n\log^2 n}\) & one pass & \coderef{c18-833d} \\
Engine pick table (which tree engine) & -- & -- & \coderef{c18-map} \\
Centroid decomposition kit (one listing) & -- & -- & \coderef{c18-kit} \\
CF 757G (Summit) -- range sum of distances & \(\bigO{n\log n}\) & \(\bigO{\log^2 n}\) & \coderef{c18-summit} \\
\bottomrule
\end{tabular}
\end{center}

\chaptersection{Master Decision Flowchart}

Post-study ladder: walk top to bottom, stop at the first branch that matches. It decides
\emph{whether the problem quantifies over all paths} (centroid's home) and, if so,
\emph{which per-centroid bookkeeping} the answer needs -- and hands back every cell that
Chapters~13/17/22/25/27 already own. A \emph{single} chosen path with range updates goes to
heavy-light (\(\chref{c16}\)); a subtree set statistic goes to sack (\(\chref{c28}\)); a
changing topology goes to link-cut (\(\chref{c23}\)).

\begin{lstlisting}[language={}, basicstyle=\ttfamily\small, numbers=none, frame=single]
  Does the problem ask about PATHS in a FIXED tree?
  |
  |-- ONE chosen path that you UPDATE and QUERY (sum/min/assign on a subpath)?
  |     +-- YES: HEAVY-LIGHT + a segtree (Ch.22, c16).  NOT centroid -- centroid
  |         cannot restrict work to one path's nodes (def:c18-map, secref:c18-boundary).
  |
  |-- A SUBTREE (everything under some node), usually a SET statistic (distinct
  |   colours, k-th largest)?
  |     +-- YES: DSU-ON-TREE / sack (Ch.17, c28) or Euler-tour + Fenwick (Ch.18, c14).
  |         Sack MERGES children bottom-up; centroid SPLITS at the centre top-down.
  |
  |-- A STATE PER NODE from a LOCAL recurrence (independent set, matching)?
  |     +-- YES: TREE DP (Ch.13, c10); change the root -> REROOTING (Ch.20, c20).
  |         No distance / all-pairs dimension, so halving buys nothing.
  |
  |-- ALL simple paths (or a distance-to-a-SET), with a value ADDITIVE along the
  |   path -- THIS chapter.  First BUILD the centroid tree (24.1):
  |     decompose(entry,pc): tot=size; c=find_centroid (Ch.12 thm:c09-centfind);
  |     par_c[c]=pc; gather dist(x,c) into x's O(log n) ancestor tables; dead[c]=1;
  |     recurse on each surviving piece.  lst:c18-build.
  |
  |     |-- COUNT / SUM over pairs by DISTANCE?  (24.2)
  |     |     At each centroid claim the paths THROUGH it: build one arm's depth
  |     |     histogram, QUERY-then-MERGE so same-arm pairs are never formed here,
  |     |     then SUBTRACT each arm re-counted.  CF 161D (=K) is Ch.13's DP --
  |     |     the centroid upgrade is prose only.  lst:c18-count, lst:c18-715c.
  |     |
  |     |-- DYNAMIC nearest/farthest marked node, online, FIXED tree?  (24.3)
  |     |     One distance table per centroid; mark/query walk a node's O(log n)
  |     |     centroid ancestors taking the min.  CF 342E Xenia is Ch.12's URL.
  |     |     lst:c18-index.
  |     |
  |     |-- Pair count with a TWO-sided budget (<= L edges AND <= W weight)?  (24.3)
  |     |     Sort the (edges,weight) table on edges, sweep, Fenwick over compressed
  |     |     weight for the partner count; ordered minus self, then halve.  CF 293E.
  |     |     lst:c18-293e.
  |     |
  |     |-- Aggregate over all paths a SINGLE TAG cannot fold (product, ratio)?  (24.4)
  |     |     General plug-in template: FEATURE (additive per node), ACCEPT (2-D
  |     |     dominance), FOLD (accumulator).  CF 833D folds a product via modular
  |     |     division for the same-arm subtraction.  lst:c18-beyond, lst:c18-833d.
  |     |
  |     |-- SUM of dist(a_i, v) over a POSITION RANGE of a permutation, online?
  |            -> SUMMIT CF 757G: Fenwick per centroid over POSITIONS + the child-
  |               "sub" subtraction (all - sub[child]).  lst:c18-summit.
  |
  |-- NONE of the above -- you probably do NOT need centroid decomposition:
        a per-vertex 0/1 knapsack over a tree is NOT all-pairs; a MOVING topology is
        link-cut (Ch.25, c23); an OFFLINE batch of node-set queries is Mo (Ch.27, c21).
\end{lstlisting}

\chaptersection{Common Mistakes Checklist}

\begin{itemize}[nosep, leftmargin=1.4em]
  \item \textbf{Reusing subtree sizes across levels.} The size array is only valid inside the
        \emph{current} undead component; recompute \(\mathit{dfs\_size}\) at every
        \(\mathit{decompose}\) call with a fresh entry. A stale \(\mathit{sz}\) picks the wrong
        centroid and the height bound \(\thmref{c18-height}\) collapses to a chain. \(\coderef{c18-build}\).
  \item \textbf{Forgetting the \code{dead[]} marker.} Once a centroid is claimed it must be
        skipped by every downstream DFS. Enter a removed node and the arm collection walks
        \emph{up} into an already-processed component -- double counting or infinite recursion.
        \(\defref{c18-decomp}\), \(\coderef{c18-build}\).
  \item \textbf{Missing the same-arm subtraction.} At a centroid, a pair whose two endpoints
        live in the \emph{same} child arm is not a path through this centroid -- it belongs to a
        lower one. Count the merged histogram and \emph{subtract} each arm re-counted, or every
        same-arm pair is claimed twice. \(\defref{c18-claim}\), \(\coderef{c18-count}\).
  \item \textbf{The parent-side size \(n-\mathit{sz}[u]\) bug.} The centroid-finder walks into
        the child with \(\mathit{sz}> \mathit{tot}/2\), but the component ``above'' a candidate
        has size \(\mathit{tot}-\mathit{sz}[u]\), not \(\mathit{sz}[u]\). Ignore it and you accept
        a node that is not a centroid. Use \(\thmref{c09-centfind}\) exactly. \(\coderef{c18-build}\).
  \item \textbf{Query-then-merge order flipped.} In the histogram count you must \emph{query}
        the accumulated arms for a new arm's depths \emph{before} merging that arm in; merge
        first and a pair from the same arm matches against itself. \(\coderef{c18-count}\).
  \item \textbf{Indexing the distance table with an array \emph{position} that is not the
        node id.} The centroid tables are keyed by node; the Summit's Fenwick is keyed by
        \emph{array position} and a node moves between positions on a swap. Conflating the two
        silently reads the wrong bucket. \(\coderef{c18-summit}\).
  \item \textbf{Not dividing the ordered count by two (or forgetting the self-match).} A
        two-pointer / dominance sweep counts \((u,v)\) and \((v,u)\); the diagonal \(u=v\) can
        also sneak in when \(2a\le L\) and \(2b\le W\). Subtract self first, then halve. \(\coderef{c18-293e}\).
  \item \textbf{Product fold with integer division.} The same-arm correction on a
        \emph{multiplicative} accumulator is division modulo the prime -- use a modular inverse,
        not \code{/}. Integer division truncates and destroys the residue. \(\coderef{c18-833d}\).
  \item \textbf{Off-by-one in the centroid-depth badge.} CF 321C prints \code{'A'+depth}; if the
        top centroid is depth \(0\) the root prints \(\mathtt{A}\). Start depth at \(0\), not
        \(1\), or every letter shifts. \(\coderef{c18-321c}\).
\end{itemize}

\chaptersection{WA Clinic}

Verdicts are symptoms; the causes are the short list this chapter drilled. Diagnose top-down --
every leaf names a listing.

\begin{figure}[H]
\centering
\begin{tikzpicture}[diagram,
    v/.style={boxnode, fill=primaryblue!12, draw=primaryblue, text width=22mm, align=center},
    b/.style={boxnode, fill=difforange!12, draw=difforange, text width=12mm},
    l/.style={boxnode, fill=successgreen!10, draw=successgreen, text width=48mm,
              font=\scriptsize, align=left}]
  \node[v] (verdict) {\textbf{Verdict on the judge}};
  \node[b, below left=8mm and 5mm of verdict] (wa) {\textbf{WA}};
  \node[b, below right=8mm and 5mm of verdict] (tle) {\textbf{TLE / RE}};
  \node[l, below right=1mm and 2mm of wa.south] (wa1)
    {the same-arm pairs were never subtracted -- a pair inside one child arm is claimed at this
     centroid and again below (\(\defref{c18-claim}\), \(\coderef{c18-count}\))};
  \node[l, below left=1mm and 2mm of wa.south] (wa2)
    {an index written under the array position instead of the node id, or the ordered count
     halved before the self-match was removed (\(\coderef{c18-index}\), \(\coderef{c18-293e}\))};
  \node[l, below right=1mm and 2mm of tle.south] (tle1)
    {a stale size array picked a non-centroid: levels degrade from \(\log n\) to \(n\), turning
     \(O(n\log n)\) into \(O(n^2)\) (\(\thmref{c18-height}\), \(\coderef{c18-build}\))};
  \node[l, below left=1mm and 2mm of tle.south] (tle2)
    {reusing \code{dead[]} DFS without recomputing per piece, or a product fold using integer
     division instead of a modular inverse (\(\coderef{c18-build}\), \(\coderef{c18-833d}\))};
  \draw[flow] (verdict) -- (wa);
  \draw[flow] (verdict) -- (tle);
  \draw[flow] (wa) -- (wa1);
  \draw[flow] (wa) -- (wa2);
  \draw[flow] (tle) -- (tle1);
  \draw[flow] (tle) -- (tle2);
\end{tikzpicture}
\caption*{WA: the same-arm correction is missing, or a distance table is indexed by the wrong key
(position vs node) or halved before removing the diagonal. TLE or RE: a stale size array picked a
non-centroid so the decomposition stops halving, or the dead marker is reused instead of recomputed.}
\end{figure}

\chaptersection{Quick Reference Card}

\begin{quickref}[Chapter 24 at a glance]
\textbf{When a query quantifies over ALL paths (or a distance-to-a-set) on a fixed tree, split the
tree at its centroid, claim each path at the highest centroid on it, recurse, and merge a
per-component distance table -- \(\bigO{\log n}\) levels. }

\medskip
\textbf{Build (24.1).} \(\mathit{decompose}(\mathit{entry},\mathit{pc})\): recompute the
component size, \(\mathit{find\_centroid}\) (Ch.~12's \(\thmref{c09-centfind}\) as a black box),
set \(\mathit{par\_c}\), gather distances, mark \code{dead}, recurse on each piece. Height
\(\le\lceil\log_2 n\rceil\) (\(\defref{c18-decomp}\), \(\thmref{c18-height}\), \(\coderef{c18-build}\)).

\medskip
\textbf{Count (24.2).} Claim each path at its top centroid: collect every child arm's depths,
\emph{query} the running histogram then \emph{merge}, and \emph{subtract} each arm re-counted.
\(\bigO{n\log n}\) for \(\{=K,\le K,\)\text{divisible-by-}\(m\}\)
(\(\defref{c18-claim}\), \(\coderef{c18-count}\), \(\coderef{c18-715c}\)).

\medskip
\textbf{Index (24.3).} One distance table per centroid; a node lives in its \(\bigO{\log n}\)
centroid-ancestor entries. Mark in \(\bigO{\log n}\); nearest-marked = min over ancestors
(exact, \(\thmref{c18-nn}\)); a two-sided budget becomes 2-D dominance (sort on edges, Fenwick on
weight) (\(\defref{c18-index}\), \(\coderef{c18-index}\), \(\coderef{c18-293e}\)).

\medskip
\textbf{Beyond (24.4).} The template is three plugs: \emph{feature} additive along the arm,
\emph{accept} a dominance / threshold test, \emph{fold} an accumulator (sum, gcd, product via
modular inverse). Only the plugs change (\(\defref{c18-beyond}\), \(\coderef{c18-beyond}\), \(\coderef{c18-833d}\)).

\medskip
\textbf{The boundary (24.5).} Single path range \(\to\) HLD (\(\chref{c16}\)); subtree set
\(\to\) sack (\(\chref{c28}\)); per-node state \(\to\) DP/reroot (\(\chref{c10}\), \(\chref{c20}\));
offline node-sets \(\to\) Mo (\(\chref{c21}\)); moving edges \(\to\) link-cut (\(\chref{c23}\));
\emph{all paths} / distance-to-set \(\to\) centroid (\(\defref{c18-map}\), \(\secref{c18-boundary}\)).
\end{quickref}

\begin{lstlisting}[caption={Chapter 24 kit: the centroid find and the decompose-and-recurse build, the query-then-merge arm count, the per-centroid distance index with its nearest-marked query, the general D-C-on-trees skeleton, the Summit all-minus-sub range accumulator, and a Fenwick -- one compact skeleton each, in one place.}, label={lst:c18-kit}]
// ===== CENTROID FIND + DECOMPOSE (24.1) =====
int dfsSize(int u,int p){ sz[u]=1; for(int v:adj[u]) if(v!=p&&!dead[v]) sz[u]+=dfsSize(v,u); return sz[u]; }
int findCentroid(int u,int p,int tot){
    for(int v:adj[u]) if(v!=p&&!dead[v]&&sz[v]>tot/2) return findCentroid(v,u,tot);
    return u;                                  // no child > tot/2 and tot-sz[u] <= tot/2  (Ch.12)
}
void decompose(int entry,int pc){
    int tot=dfsSize(entry,0), c=findCentroid(entry,0,tot);
    par_c[c]=pc;                               // centroid-tree parent
    dead[c]=true;                              // CLAIM c: no downstream DFS may re-enter
    for(int v:adj[c]) if(!dead[v]) decompose(v,c);
}

// ===== COUNT pairs at distance K through the top centroid (24.2) =====
ll count(int c,int K){
    vector<int> hist(K+1,0); ll res=0;
    for(int arm: children(c)){
        vector<int> d=depths(arm,c);           // distances from c into this arm
        for(int x: d) if(x<=K) res+=hist[K-x]; // QUERY accumulated arms BEFORE merge
        for(int x: d) if(x<=K) hist[x]++;      // MERGE
    }
    return res;                                // pairs THROUGH c; same-arm handled by recursion
}

// ===== DISTANCE INDEX: nearest marked (24.3) =====
void build(int c,int entry){ dfsDist(entry,-1,0,c); }   // push (c,dist(x,c)) into anc[x]
void mark(int x){ for(auto&[c,d]:anc[x]) best[c]=min(best[c],d); }
ll query(int x){ ll r=INF; for(auto&[c,d]:anc[x]) r=min(r,best[c]+d); return r; }

// ===== GENERAL D&C-ON-TREES skeleton (24.4) =====
//   FEAT arm(x) -> feature (additive along path) ; ACCEPT(p,q) -> dominance ; FOLD(acc,val)
void solveComp(int c){
    vector<Feat> all; ll acc=NEUTRAL;
    for(int arm: children(c)){ vector<Feat> A=collect(arm,c);
        for(auto&p:A) for(auto&q:all) if(ACCEPT(p,q)) acc=FOLD(acc, combine(p,q));
        insertAll(all,A); }                    // merge arms; subtract re-counts per arm
    recurse(pieces(c));
}

// ===== SUMMIT accumulator: all - sub[child], keyed by POSITION (24.4 / 757G) =====
void addNode(int pos,int u,int sgn){ for(k=0;k<anc[u].size();k++){ int c=anc[u][k].first; ll W=anc[u][k].second;
        all[c].add(pos,sgn,sgn*W);             // store dist(u,c) at array index pos
        if(k>0) sub[anc[u][k].c].add(pos,sgn,sgn*anc[u][k-1].second); } }  // store dist(u,par_c[c])
ll queryRange(int L,int R,int v){ ll res=0; int m=anc[v].size();
    for(int k=0;k<m;k++){ int c=anc[v][k].c; ll W=anc[v][k].d;
        ll cnt=all[c].cntQ(L,R), sm=all[c].sumQ(L,R);
        if(k+1<m){ int ch=anc[v][k+1].c; cnt-=sub[ch].cntQ(L,R); sm-=sub[ch].sumQ(L,R); }
        res+=cnt*W+sm; }                        // paths through c to v, minus same-child side
    return res; }

// ===== FENWICK (Ch.18 owned; carried forward) =====
struct Fenwick{ int n; vector<ll> t; Fenwick(int n):n(n),t(n+1,0){}
    void add(int i,ll d){ for(;i<=n;i+=i&-i) t[i]+=d; }
    ll pref(int i){ ll s=0; for(;i>0;i-=i&-i) s+=t[i]; return s; }
    ll range(int l,int r){ return pref(r)-pref(l-1); } };
\end{lstlisting}

\chaptersection{Write-From-Memory Drills}

Five signature skeletons, each lifted from a listing already in the chapter. Cover the marked line,
reproduce from memory, then check the answer one line below.

\begin{lstlisting}[caption={Drill 1 -- decompose-and-recurse: find the centroid, claim it, recurse on each piece (from lst:c18-build)}]
void decompose(int entry,int pc){
    int tot=dfsSize(entry,0);                    // size of THIS undead component
    int c=findCentroid(entry,0,tot);             // Ch.12 black box (thm:c09-centfind)
    par_c[c]=pc;                                 // link into the centroid tree
    dead[c]=true;                                // CLAIM c before recursing
    for(int v:adj[c]) if(!dead[v]) decompose(v,c);
}
\end{lstlisting}
\noindent{\scriptsize Key: recompute \(\mathit{dfs\_size}\) from the \emph{entry}, mark the centroid
\emph{dead} before descending, and recurse only into still-alive neighbours -- every piece then has
\(\le n/2\) nodes, giving height \(\le\lceil\log_2 n\rceil\) (\(\thmref{c18-height}\), \(\coderef{c18-build}\)).}

\begin{lstlisting}[caption={Drill 2 -- the arm histogram count: query accumulated arms, then merge (from lst:c18-count)}]
for each child arm of centroid c:
    d = depths of nodes in this arm (distance from c)
    for x in d:  res += hist[K - x]              // QUERY first: partner is in an EARLIER arm
    for x in d:  hist[x]++                       // MERGE this arm into the histogram
\end{lstlisting}
\noindent{\scriptsize Key: \emph{query-then-merge} is what makes same-arm pairs impossible at this
centroid -- a pair whose two endpoints share an arm is claimed lower down. Query after merging would
let an arm match itself (\(\defref{c18-claim}\), \(\coderef{c18-count}\)).}

\begin{lstlisting}[caption={Drill 3 -- the distance index: mark and nearest over centroid ancestors (from lst:c18-index)}]
void mark(int x){ for((c,d) in anc[x]) best[c]=min(best[c], d); }   // d = dist(x,c)
ll nearest(int x){ ll r=INF;
    for((c,d) in anc[x]) r=min(r, best[c]+d); return r; }
\end{lstlisting}
\noindent{\scriptsize Key: a node touches only its \(\bigO{\log n}\) centroid-ancestor tables. The
min over \(\mathit{best}[c]+\dist(x,c)\) is \emph{exact} because the top shared centroid of \(x\) and
its nearest marked node is one of \(x\)'s ancestors (\(\thmref{c18-nn}\), \(\coderef{c18-index}\)).}

\begin{lstlisting}[caption={Drill 4 -- 2-D dominance pair count: sort on one axis, Fenwick the other (from lst:c18-293e)}]
sort points by a ascending                        // a = edges to the centroid
Fenwick over compressed b                         // b = weight to the centroid
for each point q in sorted order:                 // sweep the EDGE budget
    partners = number of earlier points with b <= W - q.b   // Fenwick prefix on weight
    cnt += partners                               // pairs clearing BOTH budgets with q
self = (2*q.a <= L && 2*q.b <= W) ? 1 : 0         // did q match itself?
ordered = cnt - self                              // drop the diagonal FIRST
answer  = ordered / 2                             // unordered pairs: then halve
\end{lstlisting}
\noindent{\scriptsize Key: fix one budget by the sweep order, count the other with a Fenwick; subtract
the diagonal \((u=u)\) \emph{before} halving \((u,v)/(v,u)\) -- then re-count the same arm and divide
out to claim each pair once (\(\coderef{c18-293e}\)).}

\begin{lstlisting}[caption={Drill 5 -- product fold with same-arm correction (from lst:c18-833d)}]
tot = FOLD over all arms merged                 // whole component through c
for arm: tot = tot * inverse(FOLD over that arm) // modular DIVISION removes same-arm pairs
\end{lstlisting}
\noindent{\scriptsize Key: the accumulator is a \emph{product}, so subtracting a re-counted arm is not
a minus but a multiply-by-modular-inverse \(\mathit{mod}\;\mathit{MOD}\). Integer division truncates
and wrecks the residue -- this is the classic red on CF 833D (\(\defref{c18-beyond}\), \(\coderef{c18-833d}\)).}

\chaptersection{Complete Problem Checklist}

Every problem this chapter offers -- four fully solved in-section, worked entries, and one Summit.
Tick \emph{only} after an unaided solve; a hinted attempt counts as half.

\begin{center}
\footnotesize
\setlength{\tabcolsep}{4pt}
\renewcommand{\arraystretch}{1.25}
\begin{tabular}{@{}c p{0.16\linewidth} l p{0.20\linewidth} >{\raggedright\arraybackslash}p{0.26\linewidth} c@{}}
\toprule
\textbf{\#} & \textbf{Problem} & \textbf{Source} & \textbf{Key technique} & \textbf{What to look for} & \textbf{Done?} \\
\midrule
\multicolumn{6}{l}{\textbf{Section 24.1 -- building the centroid tree}} \\
W24.1 & decompose the running tree T (15 nodes) & worked & \lv{4} find, mark dead, recurse & top centroid 4; par\_c chain; depth \(0,1,2,3\le\lceil\log_2 15\rceil\) & \(\square\) \\
W24.1b & Ciel the Commander & CF 321C & \lv{4} print \code{'A'+cd\_depth} & the \(\mathit{decompose}\) build \emph{is} the solve; rank \(0=\) A & \(\square\) \\
\midrule
\multicolumn{6}{l}{\textbf{Section 24.2 -- counting paths by distance}} \\
W24.2 & distance-2 pairs on T & worked & \lv{4} histogram, query-then-merge & 17 pairs total; top centroid 4 claims 9, recursion 8 & \(\square\) \\
W24.2b & Digit Tree & CF 715C & \lv{5} decimal-hash pair count & forward/back hash, \(\gcd(M,10)=1\) inverse; short-circuit \(M=1\) & \(\square\) \\
\midrule
\multicolumn{6}{l}{\textbf{Section 24.3 -- the centroid tree as an index}} \\
W24.3 & nearest-marked walk on T & worked & \lv{4} ancestor table min & mark \(\{5,13\}\), query 6 \(\to\) 2 (nearest marked is node 5) & \(\square\) \\
W24.3b & Close Vertices & CF 293E & \lv{5} 2-D dominance pair count & \(\le l\) edges AND \(\le w\) weight; sort+Fenwick, subtract self then halve & \(\square\) \\
\midrule
\multicolumn{6}{l}{\textbf{Section 24.4 -- beyond path sums}} \\
W24.4 & colour-balanced product on T & worked & \lv{5} feature/accept/fold & max\((r,b)\le 2\min(r,b)\) = two dominance features; arm modular-inverse & \(\square\) \\
W24.4b & Red-Black Cobweb & CF 833D & \lv{5} product fold via inverse & clamminess product; sample \(\to 1265625\) & \(\square\) \\
\midrule
\multicolumn{6}{l}{\textbf{Section 24.5 -- the engine boundary (worked only)}} \\
W24.5 & place five specs on the tree-path map & worked & \lv{5} which engine & path range/scope: HLD/sack/DP/centroid/link-cut -- pick the cell & \(\square\) \\
\midrule
\multicolumn{6}{l}{\textbf{Summit (chapter end)}} \\
W24 & Can Bash Save the Day? & CF 757G & \lv{6} centroid + Fenwick, all-sub & sum dist\((a_i,v)\) over position range under swaps; XOR-decoded; \(23,37,28\) & \(\square\) \\
\bottomrule
\end{tabular}
\end{center}

\chaptersection{Recommended Study Order}

Five focused days. Day~1 internalises the centroid-tree build and CF 321C; day~2 the
query-then-merge distance count and CF 715C; day~3 the distance index and CF 293E; day~4 the
general template and CF 833D plus the engine boundary; day~5 the Summit and the mixed set.

\begin{center}
\small
\renewcommand{\arraystretch}{1.25}
\begin{tabular}{@{}llp{0.55\linewidth}@{}}
\toprule
\textbf{Day} & \textbf{Focus} & \textbf{Tasks} \\
\midrule
1 & Read + do 24.1 & decompose T by hand to \(\mathit{par\_c}\) and depths (\(\figref{c18-centtree}\)); write \(\coderef{c18-build}\) cold; then CF 321C printing \code{'A'+depth} (\(\coderef{c18-321c}\)) \\
2 & Read + do 24.2 & hand-count distance-2 pairs on T (\(17=9+8\)); name the query-then-merge and same-arm steps (\(\coderef{c18-count}\)); then CF 715C with the forward/backward decimal hash (\(\coderef{c18-715c}\)) \\
3 & Read + do 24.3 & trace the nearest-marked walk (mark \(\{5,13\}\), query 6 \(\to\) 2) (\(\coderef{c18-index}\)); write the ancestor index cold; then CF 293E's sort-plus-Fenwick dominance count (\(\coderef{c18-293e}\)) \\
4 & Read + do 24.4 + 24.5 & name the three plugs (feature/accept/fold) on CF 833D and replay sample 1 to \(1265625\) (\(\coderef{c18-833d}\)); place five specs on the engine map (\(\defref{c18-map}\)) \\
5 & Summit + mixed & derive the all-minus-sub accumulator and replay the decoded samples \(23,37,28\) (\(\probref{c18-summit}\)); then the first 3 mixed entries \\
\bottomrule
\end{tabular}
\end{center}

\chaptersection{Difficulty Progression}

Tiers are competency gates -- each row names the \emph{skill} the tier unlocks.

\begin{center}
\small
\renewcommand{\arraystretch}{1.3}
\begin{tabular}{@{}lp{0.32\linewidth}p{0.40\linewidth}@{}}
\toprule
\textbf{Tier} & \textbf{Problems} & \textbf{What passing this tier proves} \\
\midrule
\lv{3} \ \(\bigstar\bigstar\bigstar\)
  & the \(\mathit{find\_centroid}\) walk; a single arm's depth collection
  & you recompute sizes per component and pick the true balance point by hand \\
\lv{4} \ \(\bigstar\bigstar\bigstar\bigstar\)
  & CF 321C; the distance-\(K\) histogram on T; the nearest-marked trace
  & you build the centroid tree, keep the query-then-merge order, and walk ancestor tables \\
\lv{5} \ \(\bigstar\bigstar\bigstar\bigstar\bigstar\)
  & CF 715C, CF 293E, CF 833D, mixed at CF~2400--2900
  & you bolt an ordered structure (hash, Fenwick, product) onto each centroid bucket \\
\lv{6} \ \(\bigstar\bigstar\bigstar\bigstar\bigstar\bigstar\)
  & CF 757G (Summit), the CF~3100--3400 mixed set
  & you keep \emph{two} accumulators (all and sub) keyed by position and fold them correctly \\
\bottomrule
\end{tabular}
\end{center}

\chaptersection{Reduction Gallery}

Every entry is ``one all-paths shape \(\to\) one centroid move.'' Each line collapses a problem
family onto a listing already in the chapter.

\begin{itemize}[nosep, leftmargin=1.4em]
  \item ``output the centroid hierarchy / rank every node by the level that removed it'' (CF 321C)
        \(\to\) the build \emph{is} the solve: \(\mathit{decompose}\) writes \(\mathit{par\_c}\) and
        depth, print \code{'A'+depth} (\(\coderef{c18-build}\), \(\coderef{c18-321c}\)).
  \item ``count unordered pairs at distance \(=K\) / \(\le K\)'' (CF 161D, Ch.~13's DP upgraded) \(\to\)
        claim each pair at its top centroid, query-then-merge a depth histogram, subtract the same
        arm (\(\defref{c18-claim}\), \(\coderef{c18-count}\)).
  \item ``count pairs whose path-as-a-number is divisible by \(m\)'' (CF 715C) \(\to\) the same
        centroid count with a forward/backward decimal hash and \(\gcd(m,10)=1\) inverse
        (\(\coderef{c18-715c}\)).
  \item ``distance-to-a-set / nearest marked node, online, fixed tree'' (CF 342E Xenia, Ch.~12's URL)
        \(\to\) one distance table per centroid, min over a node's \(\bigO{\log n}\) ancestors
        (\(\thmref{c18-nn}\), \(\coderef{c18-index}\)).
  \item ``pairs within TWO budgets (\(\le L\) edges AND \(\le W\) weight)'' (CF 293E) \(\to\) 2-D
        dominance in the bucket: sort on edges, Fenwick over compressed weight, subtract self then
        halve (\(\coderef{c18-293e}\)).
  \item ``aggregate over all paths a single tag cannot fold (product, gcd, ratio)'' (CF 833D) \(\to\)
        feature/accept/fold plugs; the same-arm correction is a modular inverse on a product
        (\(\defref{c18-beyond}\), \(\coderef{c18-833d}\)).
  \item ``sum of distances over a \emph{position range} of a permutation, under swaps'' (CF 757G,
        Summit) \(\to\) a Fenwick per centroid keyed by position, with the child \emph{sub}
        table so each pair is claimed once (\(\coderef{c18-summit}\)).
  \item ``is this even a centroid job?'' \(\to\) check the engine map: single path range \(\to\) HLD,
        subtree set \(\to\) sack, per-node state \(\to\) DP, moving edges \(\to\) link-cut
        (\(\defref{c18-map}\), \(\secref{c18-boundary}\)).
\end{itemize}

\chaptersection{Mixed Problem Set}

Ten entries, each live-verified against the current problem sets (Codeforces tags
\emph{divide and conquer} plus \emph{trees}; CF has no dedicated centroid tag) and absent from
every other roster in the book (the owned IDs -- CF 342E in Ch.~12, CF 161D in Ch.~13, CF
914D/1601C/1743F, and this chapter's own solved 321C / 715C / 293E / 833D -- are excluded). None
says ``use centroid decomposition'' -- you recognise it from \emph{what the query quantifies over}
(all paths, or a distance-to-a-set) and \emph{whether the topology is fixed}. Each hint names its
reduction and the Chapter~24 listing it leans on, ordered ascending by effective difficulty.

\begin{mixedset}
  \practice{Fib-tree}{CF 1491E}
    {https://codeforces.com/problemset/problem/1491/E}{CF 2400 L5}
    {Cut a tree into two connected parts repeatedly so the part sizes follow a Fibonacci-like
     rule. Each cut is ``split the current component at a node,'' so the natural candidate is the
     balance split of \(\secref{c18-tree}\): the centroid is a node whose removal leaves pieces of
     size \(\le n/2\), and a greedy centroid-first cut is exactly the move a Fibonacci decomposition
     needs. Reuse \(\coderef{c18-build}\)'s recompute-sizes-per-component discipline; the difference
     is that you cut at a \emph{node} (edge removal), validating sizes, not merely ranking.}
  \practice{New Year Tree}{CF 379F}
    {https://codeforces.com/problemset/problem/379/F}{CF 2400 L5}
    {The tree only \emph{grows} -- you attach a new leaf -- and queries ask for the distance between
     two existing nodes. A growing tree is a fixed tree at every instant, so \(\dist(u,v)=
     \mathit{depth}_w[u]+\mathit{depth}_w[v]-2\mathit{depth}_w[\mathrm{lca}]\) with binary lifting is
     the cheapest route; but see why centroid decomposition is the WRONG pick here
     (\(\defref{c18-map}\)): the centroid tree would have to be rebuilt after every attachment, so a
     one-shot LCA wins. Recognise the boundary, not the engine.}
  \practice{Magic Matrix}{CF 632F}
    {https://codeforces.com/problemset/problem/632/F}{CF 2400 L5}
    {A ``magic matrix'' is defined by a nearest-neighbour property you check with a \emph{divide and
     conquer on the tree}: root arbitrarily and verify, for each node, that the closest marked node
     can be found by recursion. That is the distance-index pattern -- a per-centroid nearest
     structure (\(\thmref{c18-nn}\), \(\coderef{c18-index}\)) answers each ``who is nearest?'' check
     in \(\bigO{\log n}\). Port the \(\mathit{best}[c]+d\) min from the index listing verbatim.}
  \practice{Palindromes in a Tree}{CF 914E}
    {https://codeforces.com/problemset/problem/914/E}{CF 2400 L5}
    {For every node, count the paths \emph{through} it whose label multiset is palindrome-able (at
     most one odd colour). The path is claimed at its top centroid, and you merge per-arm colour
     parities with the query-then-merge histogram of \(\coderef{c18-count}\): parity is the
     \emph{feature}, the at-most-one-odd rule is the \emph{accept}, and the count is the
     \emph{fold}. The same-arm subtraction is the one thing that keeps it correct.}
  \practice{GCD Counting}{CF 990G}
    {https://codeforces.com/problemset/problem/990/G}{CF 2400 L5}
    {Count paths whose node-value gcd is at least \(2\), by prime factor. Fix a prime \(p\) and keep
     the component of \(p\)-divisible nodes -- inside it, count connected paths, a task that a
     centroid pass over each surviving component handles with the arm merge of
     \(\coderef{c18-count}\). Mobius inversion cleans the double-count across primes; the centroid
     machinery is the ``all paths, additive along path'' engine
     (\(\defref{c18-beyond}\)).}
  \practice{Vertex Pairs}{CF 2042E}
    {https://codeforces.com/problemset/problem/2042/E}{CF 2900 L6}
    {Count pairs \((u,v)\) whose \emph{maximum} edge on the path is at most a per-node threshold --
     a two-sided condition on the path exactly like CF 293E. Claim each pair at its top centroid,
     collect \((\mathit{maxEdge}, \mathit{depth})\) features per arm, and count with the sort-plus-Fenwick
     dominance sweep of \(\coderef{c18-293e}\). The reduction is ``within a budget over all
     pairs \(\to\) 2-D dominance in the centroid bucket.''}
  \practice{Logistical Questions}{CF 566C}
    {https://codeforces.com/problemset/problem/566/C}{CF 3000 L6}
    {Find the node minimising a weighted sum of distances to all others. That node is a
     \emph{weighted centroid} (the balance point under weights), which \(\secref{c18-tree}\) and
     Ch.~12's \(\thmref{c09-centfind}\) locate by walking to the heavy side -- here, the side whose
     weighted mass exceeds half. Binary-search the answer value using the distance-table walk of
     \(\coderef{c18-index}\) to check a candidate. Centroid \emph{is} the objective, not just the tool.}
  \practice{Tavas on the Path}{CF 536E}
    {https://codeforces.com/problemset/problem/536/E}{CF 3100 L6}
    {A fixed tree with path queries whose answers must be precomputed efficiently -- but read the
     engine map first: a \emph{single} queried path with an associative combine is a heavy-light or
     static-path problem (\(\chref{c16}\)), \emph{not} an all-paths problem. Use this entry to
     practise REJECTING centroid: the query ranges over one chosen path, so \(\defref{c18-map}\)
     hands it to Ch.~22. Contrast with \(\secref{c18-count}\).}
  \practice{Statistics on Tree}{CF 2222G}
    {https://codeforces.com/problemset/problem/2222/G}{CF 3100 L6}
    {An aggregate over all root-to-leaf paths of a tree, where each path is really the unique path to
     a query node -- i.e.\ an accumulation that a per-node DP cannot scope to \emph{every} path at
     once. Claim each path at its top centroid and fold with the feature/accept/fold template of
     \(\coderef{c18-beyond}\); the bucket aggregate is a plain sum (a composable fold), so no ordered
     structure is needed -- the win is the \(\bigO{\log n}\)-level coverage of all paths.}
  \practice{Iqea}{CF 936E}
    {https://codeforces.com/problemset/problem/936/E}{CF 3400 L6}
    {The hardest all-paths entry: aggregate a value over every simple path with a filter that is a
     dominance/threshold on an additive feature -- the exact home of centroid decomposition. Run
     \(\coderef{c18-beyond}\) with the pair of additive features that encode the constraint (as CF
     833D used \(2r-b,2b-r\)), a Fenwick over the compressed second feature inside each bucket
     (\(\coderef{c18-293e}\)), and the same-arm correction. Everything from 24.2--24.4 composes here.}
\end{mixedset}


\chaptersection{The Summit -- CF 757G Can Bash Save the Day?}

\begin{notebox}[Why this is the chapter top rope]
Every in-section engine here answers one kind of question about a \emph{set of paths}; the Summit
makes you answer a question about a \emph{set of array positions} whose node contents \emph{move},
and do it online. You need \emph{two} of the chapter's ideas at once: the per-centroid distance
table that lets a query walk a node's \(\bigO{\log n}\) ancestors (\(\secref{c18-index}\)), and the
top-centroid claim that lets you \emph{subtract} the child-side so each path is counted once
(\(\defref{c18-claim}\)). Bolt a Fenwick / segment tree keyed by \emph{position} onto each centroid
bucket (Chapter~18's \(\chref{c14}\) carried forward) and hold two accumulators per centroid -- the
\texttt{all} table and the child \texttt{sub} table. Adjacent swaps are then just delete-and-reinsert
a node's \(\bigO{\log n}\) entries. It is the honest summit rather than the highest badge: it is the
hardest problem this chapter \emph{fully derives} and hand-replays to the judge's output.
\end{notebox}

\solvedproblem{c18-summit}{Can Bash Save the Day?}{CF 757G}
  {https://codeforces.com/problemset/problem/757/G}{CF 3400 L6}
  {Centroid decomposition with a Fenwick or segment tree per centroid over array positions storing each node's distance to the centroid, plus a child-side subtraction table; a range-sum-of-distances query walks the query node's O(log n) centroid ancestors combining all minus sub, and an adjacent swap re-inserts two nodes' ancestor entries; queries are XOR-encoded so the structure stays online}

\begin{problemstatement}[CF 757G Can Bash Save the Day?]
You are given a weighted tree on \(n\) nodes (edge lengths are positive integers up to
\(10^6\)) and a permutation \(a_1,\dots,a_n\) of \(1,\dots,n\). Let \(\dist(u,v)\) be the total
edge weight on the unique path between \(u\) and \(v\). Answer \(q\) online queries of two types:

\textbf{Type 1,} written \texttt{1 l r v}: report \(\sum_{i=l}^{r} \dist(a_i, v)\) -- the sum of
distances from node \(v\) to the nodes occupying array positions \(l, l+1, \dots, r\).\\
\textbf{Type 2,} written \texttt{2 x}: swap \(a_x\) and \(a_{x+1}\); the new permutation is used by
all later queries.

The queries are \textbf{XOR-encoded} against the previous answer (so you cannot reorder them). Read
the raw integers \(a,b,c\) for a type-1 query and \(a\) for a type-2 query, and decode with the
running answer \(\mathit{last}\) (initialised to \(0\), and \emph{unchanged} by a type-2 query):
\[
  l = a \oplus \mathit{last},\quad r = b \oplus \mathit{last},\quad v = c \oplus \mathit{last}
  \qquad(\text{type }1),\qquad x = a \oplus \mathit{last}\qquad(\text{type }2),
\]
where \(\oplus\) is bitwise exclusive-or. The guarantees hold after decoding:
\(1\le l\le r\le n\), \(1\le v\le n\), \(1\le x\le n-1\).

\textbf{Input.} Line~1: \(n, q\) (\(1\le n,q\le 2\cdot10^5\)). Line~2: the permutation
\(a_1,\dots,a_n\). Next \(n-1\) lines: \(u, v, w\), an edge of length \(w\) (\(1\le w\le 10^6\)).
Then \(q\) blocks: a line \texttt{t}, and a line of the (encoded) query parameters.

\textbf{Output.} For each type-1 query, print the sum on its own line.

\textbf{Sample.} \(n{=}5,q{=}5\), \(a=(4,5,1,3,2)\), edges \((4,2,4),(1,3,9),(4,1,4),(4,5,2)\);
after XOR-decoding the queries are \(\{1\ 1\ 5\ 4\}, \{1\ 1\ 3\ 3\}, \{2\ 3\}, \{2\ 2\},
\{1\ 1\ 3\ 3\}\), and the three type-1 answers are \texttt{23}, \texttt{37}, \texttt{28}.
\end{problemstatement}

\begin{attemptgate}[Attempt it yourself first -- \lv{6}, budget a real afternoon]
Summing \(\dist(a_i,v)\) one position at a time is \(\bigO{n}\) per query, i.e.\ \(\bigO{nq}\)
overall -- \tle at \(2\cdot10^5\), and each \(\dist\) would itself need an LCA. The XOR encoding
means you cannot batch offline.
\peek{a distance \(\dist(u,v)\) splits cleanly once you fix a centroid that lies \emph{on} the path
\(u\)--\(v\) -- which centroid is that, and what does \(\dist(u,v)\) become?}
\nudge{Take \(c\) to be the top shared centroid of \(u\) and \(v\): then
\(\dist(u,v)=\dist(u,c)+\dist(v,c)\). So \(\sum_{i\in[l,r]}\dist(a_i,v)=\sum_c\big(\dist(v,c)\cdot
\#\{i:[l,r],\,c\text{ top for }a_i\}+\sum \dist(a_i,c)\big)\) -- but you must claim each \(a_i\) at
\emph{exactly one} centroid \(c\).}
\fullhint{For each centroid \(c\) keep two structures keyed by array \emph{position}: \texttt{all[c]}
holds \(\dist(a_i,c)\) and a count for every position whose node lies in \(c\)'s component, and
\texttt{sub[c]} holds \(\dist(a_i,\mathit{par\_c}[c])\) for those same positions. To answer a query
for \(v\), walk \(v\)'s centroid ancestors \(c\) (root toward \(v\)); at level \(c\) the child \(d\)
of \(c\) on the path to \(v\) owns exactly the positions whose top shared centroid with \(v\) is
\emph{below} \(c\), so subtract \texttt{sub[d]}: the contribution is
\((\mathit{cnt}_{\mathit{all}[c]}-\mathit{cnt}_{\mathit{sub}[d]})\cdot\dist(v,c)+
(\mathit{sum}_{\mathit{all}[c]}-\mathit{sum}_{\mathit{sub}[d]})\). Each position lives in
\(\bigO{\log n}\) centroid tables, so a swap re-inserts two nodes in \(\bigO{\log^2 n}\) and a query
is \(\bigO{\log^2 n}\). A per-centroid Fenwick or dynamic segment tree over positions supports the
point updates and range reads; use 64-bit -- sums reach \(O(n\cdot 10^{11})\).}
\end{attemptgate}

\solstep{Understanding}
Fix a query \((l,r,v)\). For each position \(i\in[l,r]\) we must add \(\dist(a_i,v)\). The key
identity (\(\defref{c18-claim}\)) is that every node pair has a \emph{unique top shared centroid}
\(c\) lying on the path between them, and for that \(c\), \(\dist(a_i,v)=\dist(a_i,c)+\dist(v,c)\).
Group the range's positions by that top centroid \(c\). As \(c\) ranges over \(v\)'s own centroid
ancestors, the positions whose top shared centroid with \(v\) is exactly \(c\) are: those in \(c\)'s
component \emph{minus} those in the child-of-\(c\) component that also contains \(v\). For those
positions the distance-to-\(c\) values are exactly what a table at \(c\) stores, and \(\dist(v,c)\) is
a constant of the walk. So the whole sum is a difference of two per-centroid aggregates: the
\texttt{all[c]} table minus the \texttt{sub[d]} table of the child \(d\) toward \(v\). Positions move
only by adjacent swaps, which are just delete-and-reinsert on \(\bigO{\log n}\) centroid tables per
node.

\solstep{Key Observations}
\begin{observation}
Each node \(x\) is stored in exactly its \(\bigO{\log n}\) centroid-ancestor tables (the identity
behind \(\thmref{c18-nn}\)): for every centroid \(c\) on \(x\)'s path to the centroid-tree root,
\(x\in\texttt{all[c]}\) with value \(\dist(x,c)\), and \(x\in\texttt{sub[c]}\) with value
\(\dist(x,\mathit{par\_c}[c])\). That bound is what makes both updates and queries polylogarithmic.
\end{observation}
\begin{observation}
When walking \(v\)'s ancestors root\(\to v\), the ``child toward \(v\)'' at level \(c\) is simply the
\emph{next} centroid in the walk. Subtracting \(\texttt{sub}[\text{next}]\) removes precisely the
positions that will be (and must be) claimed lower down, so no pair is double-counted -- this is the
same-arm correction of \(\secref{c18-count}\), re-derived for a position range.
\end{observation}
\begin{observation}
The array index (position) is the query's range axis, but the tree node is the membership axis. The
two are linked by the permutation \(a\), which a swap changes at two positions -- so a node's
contribution must be \emph{moved} between positions, not between centroid tables (whose contents are
topological and never change).
\end{observation}

\solstep{Think It With Me}
\begin{thinkbox}
\deadend{Answer each type-1 query by looping \(i=l..r\) and summing \(\dist(a_i,v)\) via LCA}
{That is \(\bigO{n}\) per query, \(\bigO{nq}\) total -- \tle at \(2\cdot10^5\). A single \(\dist\) is
cheap, but the \emph{range sum} of \(\bigO{n}\) of them is the wall. You need the range aggregate to
be answerable without touching every position.}
\deadend{One global Fenwick over positions storing \(\dist(a_i,v)\)}
{The stored value would have to depend on the queried \(v\), which changes every query -- you cannot
precompute \(\dist(a_i,v)\) for a \(v\) you have not seen. The distance must be \emph{decomposed}
into per-centroid pieces indexed only by position.}
\deadend{Centroid \texttt{all[c]} tables but no \texttt{sub} table}
{Summing \(\texttt{all[c]}\) over \(v\)'s ancestors counts a position \(i\) once for \emph{every}
centroid that is an ancestor of both \(a_i\) and \(v\), not once at their top shared centroid --
massive over-count, \textbf{WA}. The child \texttt{sub[d]} subtraction is exactly what restricts the
claim to one centroid (\(\defref{c18-claim}\)).}
\aha{Keep \emph{two} position-indexed accumulators per centroid -- \texttt{all[c]} with
\(\dist(x,c)\) and \texttt{sub[c]} with \(\dist(x,\mathit{par\_c}[c])\) -- and answer a query by
walking \(v\)'s centroid ancestors combining \(\texttt{all[c]}-\texttt{sub[child]}\). Adjacent swaps
delete and re-insert a node's \(\bigO{\log n}\) entries. Two chapter engines (the distance index and
the top-centroid claim) plus the Fenwick compose into one \(\bigO{\log^2 n}\) query.}
\end{thinkbox}

\solstep{The Approach Ladder}
\begin{center}\renewcommand{\arraystretch}{1.25}
\begin{tabular}{@{}c p{0.24\linewidth} p{0.30\linewidth} l c@{}}
\toprule
\textbf{\#} & \textbf{Approach} & \textbf{Idea} & \textbf{Time} & \textbf{Verdict}\\
\midrule
1 & \lv{2}\ loop the range, LCA each & sum \(\dist(a_i,v)\) position by position & \(\bigO{nq\log n}\) & \tle \\
2 & \lv{4}\ global position Fenwick & one table of distances & \(\bigO{\log n}\) & \textcolor{warningred}{\ding{55}\ \textbf{WA}} \\
3 & \lv{5}\ centroid \texttt{all} only & per-centroid tables, no child subtraction & \(\bigO{\log^2 n}\) & \textcolor{warningred}{\ding{55}\ \textbf{WA}} \\
4 & \lv{6}\ \textbf{centroid all \(−\) sub, by position} & walk \(v\)'s ancestors, \(\texttt{all[c]}-\texttt{sub[child]}\) & \(\bigO{\log^2 n}\) & \ac \\
\bottomrule
\end{tabular}\end{center}
Rung~4 is the Summit: rung~3 has the right tables but the wrong \emph{scope} -- without the child
\(\texttt{sub}\) subtraction a position is claimed at several centroids at once, and the over-count
is the whole error.

\solstep{Solution}
Build the centroid decomposition once (\(\coderef{c18-build}\)): as each centroid \(c\) is chosen,
gather \(\dist(x,c)\) into every node \(x\) of its component, appending \((c,\dist(x,c))\) to
\(\mathit{anc}[x]\) in root\(\to\)leaf order. Insert every position \(i\) into the structures for its
node \(a_i\): \(\texttt{all}[c]\) gets \(\dist(a_i,c)\) and \(\texttt{sub}[c]\) gets
\(\dist(a_i,\mathit{par\_c}[c])\). A type-1 query walks \(\mathit{anc}[v]\), combining
\(\texttt{all}[c]-\texttt{sub}[\text{next}]\) over the position range. A type-2 swap deletes the two
nodes at their old positions and re-inserts them at the swapped positions. Read and write 64-bit
throughout.

\solstep{Proof of Correctness}
\begin{proofbox}
\textbf{Distance decomposition.} For any nodes \(u,v\), let \(c^\star\) be their lowest common node
in the centroid tree (equivalently the first centroid whose removal separated them, or one of them).
Then \(c^\star\) lies on the path \(u\)--\(v\), so \(\dist(u,v)=\dist(u,c^\star)+\dist(v,c^\star)\)
(\(\defref{c18-claim}\)). \textbf{Single claim.} Fix \(v\). As \(c\) ranges over \(\mathit{anc}[v]\)
(root\(\to v\)) with \(d\) the next centroid toward \(v\), the set \(\{\text{nodes in }c\text{'s
component}\}\setminus\{\text{nodes in }d\text{'s component}\}\) is exactly the nodes whose \(c^\star\)
with \(v\) equals \(c\): \(u\) is claimed at \(c\) iff \(u\) shares the ancestor \(c\) with \(v\) but
does not fall into the same child branch \(d\) (else its \(c^\star\) is lower). These differences
partition all nodes, so each position \(i\) is counted \emph{once}, at its true \(c^\star\).
\textbf{Value read.} At that level the algorithm adds
\((\mathit{cnt}_{\mathit{all}[c]}-\mathit{cnt}_{\mathit{sub}[d]})\cdot\dist(v,c)+
(\mathit{sum}_{\mathit{all}[c]}-\mathit{sum}_{\mathit{sub}[d]})
=\sum_{i\in[l,r],\,c^\star(i,v)=c}(\dist(v,c)+\dist(a_i,c))
=\sum \dist(a_i,v)\) by the decomposition. \textbf{Updates.} A swap only permutes which node sits at
which position; deleting the node's \(\bigO{\log n}\) entries and re-inserting them at the new
position keeps every \(\texttt{all}/\texttt{sub}\) table equal to the current position-assignment, so
every later range read is exact. Summing over levels gives \(\sum_{i=l}^r\dist(a_i,v)\). \(\Box\)
\end{proofbox}

\solstep{Implementation}
\begin{lstlisting}[caption={CF 757G Summit: centroid decomposition; per centroid, two dynamic segment trees over POSITIONS -- all[c] stores dist(node,c) and sub[c] stores dist(node,parc[c]). A query for v walks v's centroid ancestors combining all minus the child's sub; a swap deletes and re-inserts two nodes' O(log n) entries. XOR-decode against the last answer. ASCII only}, label={lst:c18-summit}]
#include <bits/stdc++.h>
using namespace std;
typedef long long ll;
const int MAXN = 200005;
const int MAXP = 8000000;              // segment-tree pool, O(n log^2 n) nodes
int n, q;
vector<pair<int,int> > adj[MAXN];
int a[MAXN];                            // a[position] = node currently there
bool dead[MAXN]; int sz[MAXN], parc[MAXN];
vector<pair<int,ll> > anc[MAXN];        // (centroid c, dist(node,c)) in ROOT->node order
int dfsSize(int u,int p){ sz[u]=1; for(size_t k=0;k<adj[u].size();k++){ int v=adj[u][k].first;
    if(v!=p&&!dead[v]) sz[u]+=dfsSize(v,u);} return sz[u]; }
int findCent(int u,int p,int tot){ for(size_t k=0;k<adj[u].size();k++){ int v=adj[u][k].first;
    if(v!=p&&!dead[v]&&sz[v]>tot/2) return findCent(v,u,tot);} return u; }   // Ch.12 black box
void gather(int u,int p,ll d,int c){ anc[u].push_back(make_pair(c,d));
    for(size_t k=0;k<adj[u].size();k++){ int v=adj[u][k].first, w=adj[u][k].second;
        if(v!=p&&!dead[v]) gather(v,u,d+w,c);} }
void decompose(int entry,int pc){ int tot=dfsSize(entry,0); int c=findCent(entry,0,tot);
    parc[c]=pc; gather(c,0,0,c); dead[c]=true;                 // claim c, record dist(x,c)
    for(size_t k=0;k<adj[c].size();k++){ int v=adj[c][k].first; if(!dead[v]) decompose(v,c);} }
// ---- dynamic segment tree over positions, storing (count, sum of distances) ----
struct Seg{ int l,r; ll cnt,sum; } pool[MAXP]; int Npool=0;
int allR[MAXN], subR[MAXN];
int newNode(){ ++Npool; pool[Npool].l=pool[Npool].r=0; pool[Npool].cnt=pool[Npool].sum=0; return Npool; }
void upd(int &t,int lo,int hi,int pos,ll dc,ll ds){ if(!t) t=newNode();
    pool[t].cnt+=dc; pool[t].sum+=ds; if(lo==hi) return; int mid=(lo+hi)>>1;
    if(pos<=mid) upd(pool[t].l,lo,mid,pos,dc,ds); else upd(pool[t].r,mid+1,hi,pos,dc,ds); }
pair<ll,ll> qry(int t,int lo,int hi,int L,int R){ if(!t||R<lo||hi<L) return make_pair(0LL,0LL);
    if(L<=lo&&hi<=R) return make_pair(pool[t].cnt,pool[t].sum); int mid=(lo+hi)>>1;
    pair<ll,ll> A=qry(pool[t].l,lo,mid,L,R), B=qry(pool[t].r,mid+1,hi,L,R);
    return make_pair(A.first+B.first,A.second+B.second); }
void addNode(int u,int pos,int sgn){                          // insert (sgn=+1) or delete (-1)
    for(size_t k=0;k<anc[u].size();k++){ int c=anc[u][k].first; ll W=anc[u][k].second;
        upd(allR[c],1,n,pos,sgn,(ll)sgn*W);                  // all[c] : dist(u,c) at position pos
        if(k>0) upd(subR[c],1,n,pos,sgn,(ll)sgn*anc[u][k-1].second); } } // sub[c] : dist(u,parc[c])
ll queryRange(int L,int R,int v){ ll res=0; int m=(int)anc[v].size();
    for(int k=0;k<m;k++){ int c=anc[v][k].first; ll W=anc[v][k].second;   // walk v's centroid ancestors
        pair<ll,ll> AC=qry(allR[c],1,n,L,R); ll cnt=AC.first, sm=AC.second;
        if(k+1<m){ int ch=anc[v][k+1].first;                          // child of c toward v
            pair<ll,ll> SD=qry(subR[ch],1,n,L,R); cnt-=SD.first; sm-=SD.second; }
        res+=cnt*W+sm; }                                              // claim paths at top centroid c
    return res; }
int main(){ scanf("%d%d",&n,&q); for(int i=1;i<=n;i++) scanf("%d",&a[i]);
    for(int i=1;i<n;i++){ int u,v,w; scanf("%d%d%d",&u,&v,&w);
        adj[u].push_back(make_pair(v,w)); adj[v].push_back(make_pair(u,w)); }
    decompose(1,0);
    for(int i=1;i<=n;i++) addNode(a[i],i,1);                 // seed position i with node a[i]
    ll last=0;
    for(int qi=0;qi<q;qi++){ int t; scanf("%d",&t);
        if(t==1){ ll A,B,C; scanf("%lld%lld%lld",&A,&B,&C);
            int l=(int)(A^last), r=(int)(B^last), v=(int)(C^last); if(l>r) swap(l,r);
            ll ans=queryRange(l,r,v); printf("%lld\n",ans); last=ans; }
        else { ll A; scanf("%lld",&A); int x=(int)(A^last); int u=a[x], w=a[x+1];
            addNode(u,x,-1); addNode(w,x,1); addNode(w,x+1,-1); addNode(u,x+1,1);
            a[x]=w; a[x+1]=u; } }
    return 0; }
// Two accumulators, one walk. all[c] keys dist(node,c) by ARRAY POSITION; sub[c] keys
// dist(node,parc[c]) so the child-side (claimed lower) can be subtracted. O(log n) tables per node
// => O(log^2 n) per query/swap. 64-bit: a distance sum can reach ~n*1e11, still inside long long.
\end{lstlisting}

\solstep{Trace on the Official Sample}
\begin{tracebox}[Sample: decompose the 5-node tree, then replay the three decoded type-1 queries]
Centroid of \(\{1,2,3,4,5\}\) is \(4\) (pieces \(\{2\},\{5\},\{1,3\}\)); recurse gives par\_c:
\(4\!\to\!0\), \(2,5,1\!\to\!4\), \(3\!\to\!1\). Ancestor tables (root\(\to\)leaf), each entry a
centroid and the distance to it: \(\mathit{anc}[4]=[(4,0)]\); \(\mathit{anc}[2]=[(4,4),(2,0)]\);
\(\mathit{anc}[5]=[(4,2),(5,0)]\); \(\mathit{anc}[1]=[(4,4),(1,0)]\);
\(\mathit{anc}[3]=[(4,13),(1,9),(3,0)]\). The array is \(a=(4,5,1,3,2)\) at positions \(1..5\).

\textbf{Query 1: \(l=1,r=5,v=4\).} Only level \(c=4\), \(W=0\), no child: \(\texttt{all}[4]\) spans
the whole tree so the range sum is \(\sum_i\dist(a_i,4)=0+2+4+13+4=23\), count \(5\), giving
\(5\cdot0+23=23\). \(\boxed{23}\) \(\checkmark\)

\textbf{Query 2: \(l=1,r=3,v=3\).} Positions \(1..3\) hold nodes \(4,5,1\).
Level \(c=4\), \(W=13\), child \(d=1\): \(\texttt{all}[4]\) over \([1,3]\) is count \(3\), sum
\(0+2+4=6\); \(\texttt{sub}[1]\) (nodes of component \(\{1,3\}\) stored with distance to \(4\)) over
\([1,3]\) catches node \(1\) at position \(3\): count \(1\), sum \(4\). Contribution
\((3-1)\cdot13+(6-4)=26+2=28\). Level \(c=1\), \(W=9\), child \(d=3\): \(\texttt{all}[1]\) over
\([1,3]\) is node \(1\) at position \(3\): count \(1\), sum \(\dist(1,1)=0\); \(\texttt{sub}[3]\)
(node \(3\), position \(4\)) is empty. Contribution \((1-0)\cdot9+(0-0)=9\). Level \(c=3\):
component \(\{3\}\) sits at position \(4\), outside the range, contributes \(0\). Total \(28+9=37\).
\(\boxed{37}\) \(\checkmark\)

\textbf{After the two swaps} (\(x=3\), then \(x=2\)) the array is \(a=(4,3,5,1,2)\). \textbf{Query 3:
\(l=1,r=3,v=3\)} now reads nodes \(4,3,5\): \(13+0+15=\boxed{28}\) \(\checkmark\) -- matching the
judge exactly.
\end{tracebox}

\solstep{Complexity}
\begin{complexitybox}
Decomposition is \(\bigO{n\log n}\) time (\(\thmref{c18-height}\)), storing \(\bigO{\log n}\)
ancestor entries per node. A type-1 query walks \(v\)'s \(\bigO{\log n}\) ancestors, doing two
range-reads (\(\bigO{\log n}\) each on the position segment tree) per level: \(\bigO{\log^2 n}\). A
type-2 swap re-inserts two nodes across \(\bigO{\log n}\) tables each, also \(\bigO{\log^2 n}\).
Total \(\bigO{(n+q)\log^2 n}\approx 4\cdot10^{5}\cdot 18^2 \approx 1.3\cdot10^{8}\) node touches --
AC in 5s with the small \(\log\) constant. Memory: \(\bigO{n\log^2 n}\) segment-tree nodes (the
all and sub trees) -- the reason the limit is a generous 768 MB; an implementation that shares
structure or uses per-centroid Fenwicks of touched positions is leaner. Use \texttt{long long} for
answers: a distance sum reaches \(O(n\cdot\!\max\!\dist)\approx 2\cdot10^{5}\cdot2\cdot10^{11}\) in
the theoretical worst case, though the test data keeps decoded answers inside \(2^{30}\).
\end{complexitybox}

\begin{hardeningbox}
\begin{itemize}[nosep,leftmargin=1.4em]
  \item \textbf{Never forget the \texttt{sub[child]} subtraction.} Summing \(\texttt{all[c]}\) over
        \(v\)'s ancestors claims a position at \emph{every} shared centroid, not just its top one --
        the over-count is the classic \textbf{WA}. The child table removes exactly the positions that
        belong to a lower claim (\(\defref{c18-claim}\)).
  \item \textbf{Index by POSITION, key distance by the table's centroid.} \(\texttt{all[c]}\) stores
        \(\dist(x,c)\) but is \emph{addressed} by \(x\)'s array position; \(\texttt{sub[c]}\) stores
        \(\dist(x,\mathit{par\_c}[c])\). Mixing the node id and the position -- or putting
        \(\dist(x,c)\) into \(\texttt{sub[c]}\) -- silently mis-scales every difference
        (\(\coderef{c18-summit}\)).
  \item \textbf{A swap moves a node between positions, not between components.} The centroid tables
        are topological and fixed; only \(a[\cdot]\) changes. Delete the node at its old position and
        re-insert at the new one across \emph{all} its \(\bigO{\log n}\) centroid tables, both
        accumulators -- missing the \texttt{sub} update leaves a stale distance in the parent bucket.
  \item \textbf{XOR-decode against the \emph{type-1} answer.} A type-2 query does not update
        \(\mathit{last}\); reading the raw \(a,b,c\) without decoding, or advancing \(\mathit{last}\)
        on a swap, scrambles every later query into \textbf{WA}.
  \item \textbf{64-bit everywhere, but the encoded values fit in \texttt{int}.} Parameters are
        \(<2^{30}\); the \emph{answer} needs \texttt{long long}. Compute \(\mathit{cnt}\cdot W + \mathit{sum}\)
        in 64-bit -- \(\mathit{cnt}\le 2\cdot10^5\) and a distance up to \(2\cdot10^{11}\) overflow
        \code{int} instantly.
\end{itemize}
\end{hardeningbox}

\begin{recallbox}[Chapter 24 in ten seconds]
A rooted DP is blind to paths that cross its root; centroid decomposition removes the blindness by
splitting at the balance point and claiming each path at its top centroid. Build it by
find-recurse-claim (\(\coderef{c18-build}\), height \(\le\lceil\log_2 n\rceil\), \(\thmref{c18-height}\));
count pairs at a distance by query-then-merging an arm histogram and subtracting the same arm
(\(\coderef{c18-count}\), CF 715C \(\coderef{c18-715c}\)); index a component with a per-centroid
distance table for nearest-marked and 2-D-budget queries (\(\thmref{c18-nn}\), \(\coderef{c18-index}\),
CF 293E \(\coderef{c18-293e}\)); fold anything a tag cannot -- products via modular inverse
(\(\coderef{c18-833d}\)). Single-path ranges are HLD (\(\chref{c16}\)), subtree sets are sack
(\(\chref{c28}\)), moving topologies are link-cut (\(\chref{c23}\)) (\(\defref{c18-map}\)). Summit CF
757G (\(\probref{c18-summit}\)) keeps \(\texttt{all}-\texttt{sub}\) Fenwicks keyed by position and
answers \(\sum_{i\in[l,r]}\dist(a_i,v)\) online.
\selftest{
  \item Why must \(\mathit{dfs\_size}\) be recomputed inside every \(\mathit{decompose}\) call, and
        what breaks if a stale \(\mathit{sz}\) is used? \hfill\lv{3}
  \item State the query-then-merge rule and explain why merging first double-counts. \hfill\lv{4}
  \item Why is the nearest-marked min over a node's centroid ancestors \emph{exact}? \hfill\lv{4}
  \item In CF 293E, in what order do you remove the self-match and halve, and why? \hfill\lv{5}
  \item In the Summit, what does \(\texttt{sub[c]}\) store, and why must it be keyed by position
        rather than by node id? \hfill\lv{6}
}
\end{recallbox}

\chaptersection{Tool Kit Inventory}

The machine assembled through Chapter~24. Earlier chapters' rows are carried forward; this
chapter's are appended. Pointers only.

\begin{center}
\small
\renewcommand{\arraystretch}{1.25}
\begin{tabular}{@{}p{0.23\linewidth} l c p{0.22\linewidth} >{\raggedright\arraybackslash\scriptsize}p{0.18\linewidth}@{}}
\toprule
\textbf{Tool} & \textbf{Since} & \textbf{Lv} & \textbf{Trigger phrase} & \textbf{Where} \\
\midrule
\multicolumn{5}{l}{\itshape Carried forward} \\
Centroid, existence, linear find & 12 & \lv{3} & ``the balance point this chapter recurses on'' & \thmref{c09-centfind} \\
LCA / \(\dist\) oracle & 16 & \lv{4} & ``the pairwise distance the index generalises to sets'' & \chref{c11} \\
Fenwick tree & 18 & \lv{3} & ``the bucket structure for dominance and range sums'' & \chref{c14} \\
Heavy-light decomposition & 22 & \lv{5} & ``a single path you UPDATE + QUERY, not all paths'' & \chref{c16} \\
DSU-on-tree / sack & 17 & \lv{5} & ``merge child sets of a subtree'' & \chref{c28} \\
Link-cut / Euler-tour trees & 25 & \lv{6} & ``the topology moves between operations'' & \chref{c23} \\
\midrule
\multicolumn{5}{l}{\itshape Chapter 24 -- new this chapter} \\
Decompose-and-recurse build & 24.1 & \lv{4} & ``remove the centroid, mark dead, recurse on pieces'' & \defref{c18-decomp}, \coderef{c18-build} \\
Centroid tree + log height & 24.1 & \lv{4} & ``every node touches \(\log n\) centroids'' & \thmref{c18-height} \\
Top-centroid pair count & 24.2 & \lv{4} & ``count/aggregate over all paths by distance'' & \defref{c18-claim}, \coderef{c18-count} \\
Distance-to-set index & 24.3 & \lv{4} & ``nearest / marked set, online, fixed tree'' & \thmref{c18-nn}, \coderef{c18-index} \\
2-D dominance in the bucket & 24.3 & \lv{5} & ``two budgets on the path (edges AND weight)'' & \coderef{c18-293e} \\
Feature/accept/fold template & 24.4 & \lv{5} & ``an aggregate a single tag cannot fold'' & \defref{c18-beyond}, \coderef{c18-833d} \\
The tree-path engine map & 24.5 & \lv{5} & ``which tree engine, and when to hand back'' & \defref{c18-map}, \coderef{c18-map} \\
Centroid decomposition kit (one listing) & 24 & \lv{5} & ``all the skeletons in one place'' & \coderef{c18-kit} \\
Summit: all\(-\)sub by position & Summit & \lv{6} & ``sum of distances over a position range'' & \coderef{c18-summit} \\
\bottomrule
\end{tabular}
\end{center}

\chaptersection{Related Topics}

\begin{itemize}[nosep, leftmargin=1.4em]
  \item \chref{c09} (Ch.~12, Height / Diameter / Centroid) --- backward, and the foundation: Ch.~12
        owns the centroid \emph{definition}, its existence, and the linear find
        (\(\thmref{c09-centfind}\)); \(\secref{c18-tree}\) reuses that finder as a black box and starts
        at ``remove it and recurse.'' CF 342E (Xenia and Tree) is Ch.~12's URL, named in prose only at
        \(\secref{c18-index}\).
  \item \chref{c10} (Ch.~13, Tree DP) --- backward/contrast: CF 161D (count pairs at distance \(K\))
        is Ch.~13's \(\bigO{nK}\) subtree DP; \(\secref{c18-count}\) is the centroid upgrade to
        \(\bigO{n\log n}\), re-treated in prose, never re-linked.
  \item \chref{c11} (Ch.~16, LCA) --- backward: the pairwise \(\dist(u,v)\) oracle a single query uses;
        the centroid distance index (\(\secref{c18-index}\)) generalises it to \emph{sets} of marked
        nodes.
  \item \chref{c14} (Ch.~18, Fenwick Trees) --- backward: the BIT is the per-bucket structure for the
        2-D dominance count (CF 293E, \(\coderef{c18-293e}\)) and the Summit's position-indexed
        accumulator (\(\coderef{c18-summit}\)); the BIT itself is never re-taught.
  \item \chref{c16} (Ch.~22, Heavy-Light) --- backward, and a promise \textbf{discharged}: Ch.~22's
        \(\secref{c16-boundary}\) handed path \emph{counting} and distance statistics forward to
        centroid; \(\secref{c18-count}\) and \(\secref{c18-boundary}\) pay it back. The line is the
        unit of work: HLD owns \emph{one chosen path}, centroid owns \emph{all paths}.
  \item \chref{c28} (Ch.~17, DSU on Tree / Sack) --- contrast: sack \emph{merges} child multisets
        bottom-up for a \emph{subtree} statistic; centroid \emph{dissolves} a component at its middle
        top-down. The direction of the merge is the tell (\(\defref{c18-map}\)).
  \item \chref{c29} (Ch.~23, Li Chao / beats) --- forward/contrast: a path whose aggregate is a
        \emph{line envelope} queried at a variable \(x\) is Li Chao's, though it can be hosted inside a
        centroid bucket (\(\secref{c18-boundary}\)'s last table row).
  \item \chref{c23} (Ch.~25, Link-Cut / Euler-Tour Trees) --- forward: centroid decomposition is built
        \emph{once} on a fixed topology and dies the moment edges are linked or cut; that is
        link-cut's job. \textbf{This is the Next chapter.}
  \item \chref{c21} (Ch.~27, Mo's Algorithm) --- forward: an \emph{offline} batch of queries over
        arbitrary node sets with no inverse is Mo-on-trees, not the online centroid index
        (\(\secref{c18-boundary}\)).
\end{itemize}

\begin{lastword}[Last Word --- Chapter 24]
\emph{``A rooted DP can only see the paths that cross its own root, so it is blind to every pair
that straddles it; centroid decomposition lifts that blindness by cutting the tree at its balance
point, claiming each path at the single highest centroid it passes through, solving the pieces
independently, and merging a per-component distance table -- \(\bigO{\log n}\) levels, every node
touched \(\bigO{\log n}\) times.''}

\medskip
Chapter~12 proved a centroid exists and finds one in linear time; this chapter began at ``remove it
and recurse.'' Sections 24.1--24.4 turned that single move into a family of engines: the centroid
tree and its logarithmic height (\(\thmref{c18-height}\)); the query-then-merge arm count that
counts every pair by distance (\(\secref{c18-count}\)); the per-centroid distance table that indexes
a whole component so a mark or a nearest query is \(\bigO{\log n}\) (\(\thmref{c18-nn}\)); and the
feature/accept/fold template that folds products and 2-D dominance the way 24.2 folds plain sums. The
Summit CF 757G (\(\probref{c18-summit}\)) took the index and the top-centroid claim and ran them
together over \emph{positions}, answering an online range-sum-of-distances with adjacent swaps.

\medskip
Carry three reflexes forward:
\begin{enumerate}[nosep, leftmargin=1.6em]
  \item \emph{Ask what the query ranges over.} One path \(\to\) HLD; a subtree \(\to\) sack; a
        per-node state \(\to\) DP; \emph{all paths} or a distance-to-\emph{set} \(\to\) centroid
        (\(\defref{c18-map}\)).
  \item \emph{Claim each path exactly once, at its top centroid, and subtract the same arm.} The
        histogram merge, the 2-D dominance sweep, and the Summit's \(\texttt{all}-\texttt{sub}\) are
        the \emph{same} correction wearing different hats (\(\defref{c18-claim}\)).
  \item \emph{Recompute sizes per component, and honour the dead marker.} The whole \(\bigO{\log n}\)
        bound rests on each piece being \(\le n/2\); a stale size or a re-entered centroid turns it
        into a chain (\(\coderef{c18-build}\)).
\end{enumerate}

\medskip
And three habits, drilled until boring: gather the per-centroid distance table once at build time so
every later query is a walk, not a fresh DFS; subtract the self-match \emph{before} halving an ordered
pair count; and when an accumulator is a product, the same-arm subtraction is a modular \emph{inverse},
not a division. You now own the divide-and-conquer on the tree's \emph{shape} --- the one engine
heavy-light and Li Chao could only hint at for paths that refuse to share a root.

\medskip
\textbf{Next:} \chref{c23} (Link-Cut / Euler-Tour Trees) drops the last assumption this chapter leaned
on everywhere -- that the tree is \emph{fixed}. When edges are linked and cut between operations, no
decomposition built once survives; the splay-based link-cut tree recomputes path aggregates on a
\emph{moving} forest, the dynamic companion to everything in Part~V.
\end{lastword}
